% Capítulo 2: Linealización y Diferenciales

Aproximaciones lineales de funciones diferenciables en vecindarios locales y propagación de errores mediante diferenciales.

\noindent \textit{Encuentre la diferencial de la función dada en los Ejercicios 1 a 8.}

\begin{enumerate}
    \begin{minipage}{0.45\textwidth}
    \subimport{ejercicio_01/}{ejercicio.tex}
    \subimport{ejercicio_03/}{ejercicio.tex}
    \subimport{ejercicio_05/}{ejercicio.tex}
    \subimport{ejercicio_07/}{ejercicio.tex}
    \end{minipage}%
    \begin{minipage}{0.45\textwidth}
    \setcounter{enumi}{1}
    \subimport{ejercicio_02/}{ejercicio.tex}
    \addtocounter{enumi}{1}
    \subimport{ejercicio_04/}{ejercicio.tex}
    \addtocounter{enumi}{1}
    \subimport{ejercicio_06/}{ejercicio.tex}
    \addtocounter{enumi}{1}
    \subimport{ejercicio_08/}{ejercicio.tex}
    \end{minipage}
\end{enumerate}

\noindent \textit{En los Ejercicios 9 a 16, (a) encuentre la diferencial $dy$ y (b) evalúe $dy$ para los valores dados de $x$ y $dx$.}

\begin{enumerate}
    \begin{minipage}{0.45\textwidth}
    \setcounter{enumi}{8}
    \subimport{ejercicio_09/}{ejercicio.tex}
    \addtocounter{enumi}{1}
    \subimport{ejercicio_11/}{ejercicio.tex}
    \addtocounter{enumi}{1}
    \subimport{ejercicio_13/}{ejercicio.tex}
    \addtocounter{enumi}{1}
    \subimport{ejercicio_15/}{ejercicio.tex}
    \end{minipage}%
    \begin{minipage}{0.45\textwidth}
    \setcounter{enumi}{9}
    \subimport{ejercicio_10/}{ejercicio.tex}
    \addtocounter{enumi}{1}
    \subimport{ejercicio_12/}{ejercicio.tex}
    \addtocounter{enumi}{1}
    \subimport{ejercicio_14/}{ejercicio.tex}
    \addtocounter{enumi}{1}
    \subimport{ejercicio_16/}{ejercicio.tex}
    \end{minipage}
\end{enumerate}

\noindent \textit{En los Ejercicios 17 a 20 calcule $\Delta y$ y $dy$ para los valores dados de $x$ y $dx = \Delta x$. Luego trace un diagrama como el de la Figura 2.24 para ilustrar los segmentos rectilíneos de longitudes $dx$, $dy$ y $\Delta y$.}

\begin{enumerate}
    \setcounter{enumi}{16}
    \subimport{ejercicio_17/}{ejercicio.tex}
    \subimport{ejercicio_18/}{ejercicio.tex}
    \subimport{ejercicio_19/}{ejercicio.tex}
    \subimport{ejercicio_20/}{ejercicio.tex}
\end{enumerate}

\noindent \textit{En los Ejercicios 21 y 22 calcule $\Delta y$, $dy$ y $\Delta y - dy$, para los valores dados de $x$ y para cada uno de los siguientes valores de $dx = \Delta x$: 1, 0.5, 0.1 y 0.01.}

\begin{enumerate}
    \setcounter{enumi}{20}
    \subimport{ejercicio_21/}{ejercicio.tex}
    \subimport{ejercicio_22/}{ejercicio.tex}
\end{enumerate}

\noindent \textit{En los Ejercicios 23 a 32, utilice diferenciales para encontrar un valor aproximado al número indicado.}

\begin{enumerate}
    \begin{minipage}{0.45\textwidth}
    \setcounter{enumi}{22}
    \subimport{ejercicio_23/}{ejercicio.tex}
    \addtocounter{enumi}{1}
    \subimport{ejercicio_25/}{ejercicio.tex}
    \addtocounter{enumi}{1}
    \subimport{ejercicio_27/}{ejercicio.tex}
    \addtocounter{enumi}{1}
    \subimport{ejercicio_29/}{ejercicio.tex}
    \addtocounter{enumi}{1}
    \subimport{ejercicio_31/}{ejercicio.tex}
    \end{minipage}%
    \begin{minipage}{0.45\textwidth}
    \setcounter{enumi}{23}
    \subimport{ejercicio_24/}{ejercicio.tex}
    \addtocounter{enumi}{1}
    \subimport{ejercicio_26/}{ejercicio.tex}
    \addtocounter{enumi}{1}
    \subimport{ejercicio_28/}{ejercicio.tex}
    \addtocounter{enumi}{1}
    \subimport{ejercicio_30/}{ejercicio.tex}
    \addtocounter{enumi}{1}
    \subimport{ejercicio_32/}{ejercicio.tex}
    \end{minipage}
\end{enumerate}

\noindent \textit{En los Ejercicios 33 a 36 obtenga la linealización $L(x)$ de la función dada en $x_1$.}

\begin{enumerate}
    \begin{minipage}{0.45\textwidth}
    \setcounter{enumi}{32}
    \subimport{ejercicio_33/}{ejercicio.tex}
    \addtocounter{enumi}{1}
    \subimport{ejercicio_35/}{ejercicio.tex}
    \end{minipage}%
    \begin{minipage}{0.45\textwidth}
    \setcounter{enumi}{33}
    \subimport{ejercicio_34/}{ejercicio.tex}
    \addtocounter{enumi}{1}
    \subimport{ejercicio_36/}{ejercicio.tex}
    \end{minipage}
\end{enumerate}

\noindent \textit{En los Ejercicios 37 a 40, verifique la aproximación lineal dada en $x_1 = 0$.}

\begin{enumerate}
    \begin{minipage}{0.45\textwidth}
    \setcounter{enumi}{36}
    \subimport{ejercicio_37/}{ejercicio.tex}
    \addtocounter{enumi}{1}
    \subimport{ejercicio_39/}{ejercicio.tex}
    \end{minipage}%
    \begin{minipage}{0.45\textwidth}
    \setcounter{enumi}{37}
    \subimport{ejercicio_38/}{ejercicio.tex}
    \addtocounter{enumi}{1}
    \subimport{ejercicio_40/}{ejercicio.tex}
    \end{minipage}
\end{enumerate}

\begin{enumerate}
    \setcounter{enumi}{40}
    \subimport{ejercicio_41/}{ejercicio.tex}
    \subimport{ejercicio_42/}{ejercicio.tex}
    \subimport{ejercicio_43/}{ejercicio.tex}
    \subimport{ejercicio_44/}{ejercicio.tex}
    \subimport{ejercicio_45/}{ejercicio.tex}
    \subimport{ejercicio_46/}{ejercicio.tex}
    \subimport{ejercicio_47/}{ejercicio.tex}
    \subimport{ejercicio_48/}{ejercicio.tex}
    \subimport{ejercicio_49/}{ejercicio.tex}
\end{enumerate}
